\documentclass[12pt,a4paper]{article}
\usepackage[utf8]{inputenc}
\usepackage{amsmath,amsfonts,amssymb}
\usepackage{listings}
\usepackage{xcolor}
\usepackage{geometry}
\usepackage{hyperref}
\usepackage{fancyhdr}

\geometry{margin=1in}
\pagestyle{fancy}
\fancyhf{}
\rhead{Python Tuples - Complete Guide}
\lfoot{Python for Data Science}
\rfoot{\thepage}

\lstset{
    language=Python,
    basicstyle=\ttfamily\small,
    keywordstyle=\color{blue}\bfseries,
    commentstyle=\color{green!60!black},
    stringstyle=\color{red},
    numbers=left,
    numberstyle=\tiny\color{gray},
    stepnumber=1,
    numbersep=8pt,
    showstringspaces=false,
    breaklines=true,
    frame=single,
    backgroundcolor=\color{gray!10}
}

\title{\textbf{Python Tuples - Complete Guide}}
\author{Python for Data Science - Beginner's Level}
\date{\today}

\begin{document}

\maketitle
\tableofcontents
\newpage

\section{Overview}

Tuples are \textbf{ordered}, \textbf{immutable} (unchangeable), and allow \textbf{duplicate elements}. They are similar to lists but cannot be modified after creation, making them ideal for storing data that shouldn't change.

\section{Tuple Creation}

\subsection{Basic Creation}

\begin{lstlisting}
# Empty tuple
empty_tuple = ()
empty_tuple = tuple()

# Tuple with elements
numbers = (1, 2, 3, 4, 5)
mixed_tuple = (1, "hello", 3.14, True)

# Single element tuple (comma is required!)
single_element = (42,)  # Without comma, it's just parentheses
single_element = 42,    # Alternative syntax

# Without parentheses (tuple packing)
coordinates = 10, 20, 30

# Nested tuples
nested_tuple = ((1, 2), (3, 4), (5, 6))
\end{lstlisting}

\subsection{Using tuple() Constructor}

\begin{lstlisting}
# From list
from_list = tuple([1, 2, 3, 4])  # (1, 2, 3, 4)

# From string
from_string = tuple("hello")  # ('h', 'e', 'l', 'l', 'o')

# From range
from_range = tuple(range(5))  # (0, 1, 2, 3, 4)
\end{lstlisting}

\section{Accessing Elements}

\subsection{Indexing and Slicing}

\begin{lstlisting}
fruits = ("apple", "banana", "cherry", "date")

# Positive indexing
print(fruits[0])    # "apple"
print(fruits[2])    # "cherry"

# Negative indexing
print(fruits[-1])   # "date"
print(fruits[-2])   # "cherry"

# Slicing
numbers = (0, 1, 2, 3, 4, 5, 6, 7, 8, 9)
print(numbers[2:7])     # (2, 3, 4, 5, 6)
print(numbers[::-1])    # (9, 8, 7, 6, 5, 4, 3, 2, 1, 0)
\end{lstlisting}

\section{Tuple Methods}

\subsection{Built-in Methods}

\begin{lstlisting}
fruits = ("apple", "banana", "cherry", "banana", "date")

# count() - Count occurrences of an element
banana_count = fruits.count("banana")  # 2

# index() - Find first occurrence of an element
banana_index = fruits.index("banana")  # 1

# index() with start parameter
banana_index_after = fruits.index("banana", 2)  # 3
\end{lstlisting}

\subsection{Membership Testing}

\begin{lstlisting}
fruits = ("apple", "banana", "cherry")

# in operator
print("apple" in fruits)     # True
print("grape" in fruits)     # False

# not in operator
print("grape" not in fruits) # True
\end{lstlisting}

\section{Tuple Operations}

\subsection{Mathematical Operations}

\begin{lstlisting}
# Concatenation
tuple1 = (1, 2, 3)
tuple2 = (4, 5, 6)
combined = tuple1 + tuple2  # (1, 2, 3, 4, 5, 6)

# Repetition
repeated = (1, 2) * 3  # (1, 2, 1, 2, 1, 2)

# Length
print(len((1, 2, 3, 4, 5)))  # 5

# Min, Max, Sum (for numeric tuples)
numbers = (3, 1, 4, 1, 5, 9, 2, 6)
print(min(numbers))  # 1
print(max(numbers))  # 9
print(sum(numbers))  # 31
\end{lstlisting}

\section{Tuple Unpacking}

\subsection{Basic Unpacking}

\begin{lstlisting}
# Assign tuple elements to variables
point = (10, 20)
x, y = point
print(f"x: {x}, y: {y}")  # x: 10, y: 20

# Multiple assignment
person = ("Alice", 25, "Engineer")
name, age, job = person
\end{lstlisting}

\subsection{Advanced Unpacking}

\begin{lstlisting}
# Using * for collecting multiple elements
numbers = (1, 2, 3, 4, 5)
first, *middle, last = numbers
print(f"First: {first}")    # First: 1
print(f"Middle: {middle}")  # Middle: [2, 3, 4]
print(f"Last: {last}")      # Last: 5

# Swapping variables
a, b = 10, 20
a, b = b, a  # Swap values
\end{lstlisting}

\section{Named Tuples}

\begin{lstlisting}
from collections import namedtuple

# Define a named tuple class
Point = namedtuple('Point', ['x', 'y'])
Person = namedtuple('Person', ['name', 'age', 'city'])

# Create instances
p1 = Point(10, 20)
person1 = Person("Alice", 25, "New York")

# Access by name or index
print(p1.x, p1.y)        # 10 20
print(person1.name)      # Alice

# Named tuple methods
point_dict = p1._asdict()  # Convert to dictionary
p2 = p1._replace(x=30)     # Create new with changed field
\end{lstlisting}

\section{Performance Considerations}

\begin{table}[h]
\centering
\begin{tabular}{|l|c|l|}
\hline
\textbf{Operation} & \textbf{Time Complexity} & \textbf{Notes} \\
\hline
Access by index & O(1) & Direct access \\
Search & O(n) & Linear search \\
Count & O(n) & Must check all elements \\
Concatenation & O(n+m) & Creates new tuple \\
Slicing & O(k) & k is slice length \\
\hline
\end{tabular}
\caption{Tuple Operations Time Complexity}
\end{table}

\subsection{Memory Efficiency}

Tuples are more memory efficient than lists:
\begin{itemize}
    \item Smaller memory footprint
    \item Faster creation and iteration
    \item Can be used as dictionary keys (hashable)
\end{itemize}

\section{Common Use Cases}

\subsection{Coordinates and Points}

\begin{lstlisting}
# 2D coordinates
point_2d = (10, 20)
x, y = point_2d

# 3D coordinates
point_3d = (10, 20, 30)
x, y, z = point_3d

# RGB colors
red = (255, 0, 0)
green = (0, 255, 0)
blue = (0, 0, 255)
\end{lstlisting}

\subsection{Function Return Values}

\begin{lstlisting}
def get_name_age():
    """Return multiple values as a tuple"""
    return "Alice", 25

def calculate_stats(numbers):
    """Return statistics as a tuple"""
    return min(numbers), max(numbers), sum(numbers) / len(numbers)

# Usage
name, age = get_name_age()
minimum, maximum, average = calculate_stats([1, 2, 3, 4, 5])
\end{lstlisting}

\subsection{Dictionary Keys}

\begin{lstlisting}
# Tuples can be dictionary keys (immutable)
coordinates_data = {
    (0, 0): "Origin",
    (1, 0): "Right",
    (0, 1): "Up",
    (1, 1): "Diagonal"
}

print(coordinates_data[(1, 1)])  # "Diagonal"
\end{lstlisting}

\section{Common Pitfalls}

\subsection{Single Element Tuple}

\begin{lstlisting}
# Wrong way (creates int, not tuple)
not_a_tuple = (42)
print(type(not_a_tuple))  # <class 'int'>

# Right way (comma is required)
single_tuple = (42,)
print(type(single_tuple))  # <class 'tuple'>
\end{lstlisting}

\subsection{Modifying Nested Mutable Objects}

\begin{lstlisting}
# Tuple is immutable, but contained objects might not be
tuple_with_list = ([1, 2, 3], "hello")

# This works (modifying the list inside the tuple)
tuple_with_list[0].append(4)
print(tuple_with_list)  # ([1, 2, 3, 4], 'hello')

# This doesn't work (trying to replace tuple element)
# tuple_with_list[0] = [5, 6, 7]  # TypeError
\end{lstlisting}

\section{Summary}

Tuples are essential Python data structures that provide:
\begin{itemize}
    \item \textbf{Immutability}: Cannot be changed after creation
    \item \textbf{Order}: Maintain insertion order
    \item \textbf{Hashability}: Can be used as dictionary keys
    \item \textbf{Memory efficiency}: More compact than lists
    \item \textbf{Multiple assignment}: Enable elegant unpacking
    \item \textbf{Data integrity}: Protect data from accidental modification
\end{itemize}

Use tuples when you need ordered, unchangeable collections of data!

\end{document}